%------------------------------ Ejercicio 4.59 -------------------------------
\setcounter{neje}{58}
\begin{nejer}
Sea $(X, \mathcal{A}, \mu)$ un espacio de medida localizable. Pruebe que satisface lo siguiente.
\begin{itemize}[leftmargin=1.5em]
\item[(F)]\textit{Sea $\mathcal{F}$ una colección de funciones medibles $f \colon A_f \to \mathbb{R}$, cada una definida en un conjunto medible $A_f \en \mathcal{A}$. Suponga que cualesquiera dos funciones $f_1, f_2 \en \mathcal{F}$ coinciden casi en todas partes en $A_{f_1} \cap A_{f_2}$. Entonces existe una función medible $g \colon X \to \mathbb{R}$ tal que $g|_{A_f} = f$ casi en todas partes para todo $f \en \mathcal{F}$.}
\end{itemize}
{
% \emph{Sugerencia.} Sea $\mathcal{F}$ una colección de funciones como en (F). Para $a \en \mathbb{R}$ y $f \en \mathcal{F}$ defina
% $$
% A_f^a \coloneq \{x \en A_f \mid f(x) < a\}.
% $$
% Para $q \en \mathbb{Q}$, sea $H^q \en \mathcal{A}$ una envolvente de la colección
% $
% \mathcal{E}_q \coloneq \{A_f^q \mid f \en \mathcal{F}\}.
% $
% Defina, para $a\en\mathbb{R}$, los conjuntos medibles
% $$
% X_a \coloneq \bigcup_{q \en \mathbb{Q}, \, q < a} H^q
% $$
% Pruebe lo siguiente.
% \begin{enumerate}[label = (\roman*), leftmargin=1.8em]
%     \item Si $a < b$, entonces $X_a \subseteq X_b$.
%     \item Para todo $a \en \mathbb{R}$ el conjunto medible $X_a$ es una envolvente de la colección $\mathcal{E}_a \coloneq \{A_f^a \mid f \en \mathcal{F}\}$. Por lo tanto $\mu(A_f^a \setminus X_a) = 0$ para todo $f \en \mathcal{F}$ y, si $G \en \mathcal{A}$, entonces
%     $$
%     \mu(A_f^a \setminus G) = 0 \text{ para todo } f \en \mathcal{F} \implies \mu(X_a \setminus G) = 0.
%     $$
%     \item Si $a \en \mathbb{R}$ y $E \en \mathcal{A}$, entonces
%     $$
%     \mu(A_f^a \cap E) = 0 \text{ para todo } f \en \mathcal{F} \implies \mu(X_a \cap E) = 0.
%     $$
%     \item $\mu(X_a \cap A_f \setminus A_f^a) = 0$ para todo $f \en \mathcal{F}$ y todo $a \en \mathbb{R}$.
%     \item Defina
%     $$
%     E_0 \coloneq \left( \bigcap_{r \en \mathbb{R}} X_r \right) \cup \left( X \setminus \bigcup_{s \en \mathbb{R}} X_s \right).
%     $$
%     Entonces $E_0$ es medible y $\mu(A_f \cap E_0) = 0$ para todo $f \en \mathcal{F}$.
%     \item Para $f \en \mathcal{F}$ defina el conjunto medible $E_f \subseteq A_f$ mediante
%     $$
%     E_f \coloneq (A_f \cap E_0) \cup \bigcup_{q \en \mathbb{Q}} (A_f^q \setminus X_q) \cup \bigcup_{q \en \mathbb{Q}} (X_q \cap A_f \setminus A_f^q).
%     $$
%     Entonces $\mu(E_f) = 0$.
%     \item Defina $g \colon X \to \mathbb{R}$ mediante
%     $$
%     g(x) \coloneq \begin{cases}
%     0 & \text{si } x \en E_0, \\
%     a & \text{si } x \en X_s \text{ para todo } s > a \text{ y } x \notin X_r \text{ para todo } r < a.
%     \end{cases}
%     $$
%     Entonces $g$ está bien definida y es medible, y $g = f$ en $A_f \setminus E_f$ para todo $f \en \mathcal{F}$.
% \end{enumerate}
}
\end{nejer}
\begin{sol}
Sea $\mathcal{F}$ tal que satisface las hipótesis de (F). Entonces definimos, para cada $a\en\mathbb{R}$ y $f\en\mathcal{F}$:
\[
A^a_f \coloneq \set*{x\en A_f \given f(x) < a},
\qquad \text{(también denotado $\set{f<a}$)}
\]
Como $f$ es medible, $A^a_f\en\mathcal{A}$. Luego, por localizabilidad, para cada $a\en\mathbb{R}$ existe $H^a \en\mathcal{A}$ tal que es una envolvente de la familia $\mathcal{E}^a \coloneq \set[\big]{A^a_f \given f\en\mathcal{F} }$. Para "forzar" a que estas envolventes sean monótonas, definimos para cada $a\en\mathbb{R}$:
\[
X^a \coloneq \bigcup_{q\en\mathbb{Q}_{a}} H^q,
\quad\text{donde denotamos}\quad
\mathbb{Q}_a \coloneq \set{q\en\mathbb{Q} \given q<a}
\]
Notemos que $X^a\en\mathcal{A}$ por ser unión contable de medibles; además, $a \leq b$ implica $X^a \subseteq X^b$.%
\footnote{
    En efecto, para todo $q\en\mathbb{Q}_a$ se tiene $q<b$, así que $H^q \subseteq X^b$. Al unir para $q\en\mathbb{Q}_a$ se sigue $X^a\subseteq X^b$.
}  
Fijemos $a\en\mathbb{R}$ y probemos que $X^a$ es envolvente de $\mathcal{E}^a$:
\begin{itemize}[leftmargin=1em]
\item 
Sea $f\en\mathcal{F}$ arbitraria. Entonces:
\[
A^a_f \setminus X^a 
\;=\; A^a_f\setminus\bigcup_{q\en\mathbb{Q}_a}H^q
\;=\; \bigcap_{q\en\mathbb{Q}_a} A^a_f\setminus H^q
\]
No obstante, como $A^a_f = \set{f<a}$, para cada $q\en\mathbb{Q}_a$ tenemos $A^a_f = A^q_f \cup \set{q \leq f < a}$. Así que
\begin{align*}
A^q_f \setminus H^q
&\;=\; \big(A^q_f \cup \set{q \leq f < a}\big) \setminus H^q
\;\subseteq\; \left(A^q_f\setminus H^q\right) \cup \set{q \leq f < a}\\
&\;\subseteq\; \bigg(\bigcup_{q'\en\mathbb{Q}_a} A^{q'}_f\setminus H^{q'}\bigg) \cup \set{q \leq f < a}
\end{align*}
Por lo tanto, al tomar intersección (aprovechando la distributividad: $\bigcap_x(M\cup N_x) = M \cup \bigcap_x N_x$).
\begin{equation}\label{eq:env1}
\begin{split}
A^a_f \setminus X^a 
&\;=\; \bigcap_{q\en\mathbb{Q}_a} A^a_f\setminus H^q
\;\subseteq\; \bigcap_{q\en\mathbb{Q}_a} \bigg(\bigcup_{q'\en\mathbb{Q}_a} A^{q'}_f\setminus H^{q'}\bigg) \cup \set{q \leq f < a}\\
&\;=\; \bigg(\bigcup_{q'\en\mathbb{Q}_a} A^{q'}_f\setminus H^{q'}\bigg) \cup \bigg(\bigcap_{q\en\mathbb{Q}_a}\set{q \leq f < a}\bigg)
\end{split}
\end{equation}
Pero, gracias a la densidad de $\mathbb{Q}$, $\bigcap_{q\en\mathbb{Q}_a}[q,a) = \varnothing$, por lo cual:
\[
\bigcap_{q\en\mathbb{Q}_a}\set{q \leq f < a} 
\;=\; \bigcap_{q\en\mathbb{Q}_a} f^{-1}\bigl([q,a)\bigr) 
\;=\; f^{-1}\biggl(\bigcap_{q\en\mathbb{Q}_a}[q,a)\biggr) 
\;=\; f^{-1}(\varnothing)
\;=\; \varnothing
\] 
así que la contención (\ref{eq:env1}) se convierte en
\begin{equation}\label{eq:env2}
A^a_f \setminus X^a 
\;\subseteq\; \bigcup_{q'\en\mathbb{Q}_a} A^{q'}_f\setminus H^{q'}
\end{equation}
Como $H^{q'}$ es envolvente de $\mathcal{E}^{q'} \pert A^{q'}_f$, entonces $\mu(A^{q'}_f\setminus H^{q'}) = 0$. Se sigue que el conjunto derecho en (\ref{eq:env2}) es $\mu$-nulo por ser unión contable de conjuntos $\mu$-nulos, de donde $\mu(A^a_f \setminus X^a) = 0$.

\item
Sea $G\en\mathcal{A}$ tal que $\mu(A^a_f \setminus G) = 0$ para todo $f\en\mathcal{F}$. Luego, para todos $q\en\mathbb{Q}_a$ y $f\en\mathcal{F}$, tenemos $\mu(A^q_f\setminus G) \leq \mu(A^a_f\setminus G) = 0$ (ya que $A^q_f \subseteq A^a_f$). Por lo tanto, al ser $H^q$ envolvente de $\mathcal{E}^q$, se sigue que $\mu(H^q \setminus G) = 0$, para todo $q\en\mathbb{Q}_a$. En consecuencia, el conjunto
\[
X^a \setminus G 
\;=\; \biggl(\bigcup_{q\en\mathbb{Q}_a}H^q \biggr) \setminus G
\;=\; \bigcup_{q\en\mathbb{Q}_a}H^q\setminus G
\]
es $\mu$-nulo por ser union contable de conjuntos $\mu$-nulos. Es decir, $\mu(X^a \setminus G) = 0$.
\end{itemize}
Hemos probado que $X^a$ es envolvente de $\mathcal{E}^a$. Ahora bien, definimos la función $\tilde{g} \colon X \to \ubar{\mathbb{R}}$ tal que
\[
\tilde{g}(x) \coloneq \inf\set[\big]{a\en\mathbb{R} \given x \en X^a}
\]
la cual está bien definida gracias a que todo subconjunto de $\ubar{\mathbb{R}}$ tiene un ínfimo (\cref{nejer:1.14}) recordando $\inf\varnothing = \infty$ (por vacuidad todo real extendido es cota inferior de $\varnothing$, así que $\infty$ es su mayor cota inferior).\vspace{1em}\\
\emph{Veamos que $\tilde{g}$ es medible}:\\
Por el teorema 1.21 (\cite[\S 1.2]{Salamon}), basta con probar que $\set{\tilde{g} < b}\en\mathcal{A}$ para cada $b\en\mathbb{R}$. En efecto, veremos que $\set{\tilde{g}<b} = X^b \en \mathcal{A}$:
\begin{itemize}[leftmargin=1.8em]
\item[($\subseteq$)]Si $x\en X$ es tal que $\tilde{g}(x) < b$ entonces, por definición de ínfimo, $b$ no puede ser cota inferior de $\set{a\en\mathbb{R} \given x \en X^a}$, así que existe $a<b$ tal que $x \en X^a$. Como $X^a \subseteq X^b$, se sigue $x\en X^b$.
\item[($\supseteq$)]Si $x\en X^b$, entonces por su definición existe $q\en\mathbb{Q}_b$ tal que $x\en H^q$. Como $H^q \subseteq X^q$ tenemos que $x\en X^q$, y así $\tilde{g}(x) = \inf\set{a\en\mathbb{R} \given x \en X^a} \leq q < b$.
\end{itemize}
Esto prueba que $\set{\tilde{g} < b} = X^b$, así que $\tilde{g}$ es medible.\vspace{1em}\\
\emph{Veamos que $\tilde{g}|_{A_f} \ctpeq f$ para cada $f\en\mathcal{F}$}:\\
Fijemos una $f\en\mathcal{F}$. Por lo que ya vimos, $\set*{\tilde{g}|_{A_f} < b} = X^b \cap A_f$, y por definición $\set{f < b} = A^b_f$, así que nos bastará probar que $X^b \cap A_f \ctpeq A^b_f$.%
\footnote{
    Dados $A,B\en\mathcal{A}$, se dice que $A \ctpeq B$ cuando $\mu(A \Delta B) = 0$. Notemos que, en general, se puede decir que $A \subseteq_{\text{a.e.}} B$ cuando $\mu(A \setminus B) = 0$, y así, $A \ctpeq B$ si y solo si $A \subseteq_{\text{a.e.}} B$ y $B \subseteq_{\text{a.e.}}A$ (ya que $A \Delta B = A \setminus B \,\cup\, B\setminus A$). 
    % Se puede probar de forma sencilla que la relación "$\ctpeq$" es de equivalencia en $\mathcal{A}$ y las operaciones finitas de conjuntos están bien definidas, haciendo de $\mathcal{A}/\ctpeq$ un retículo booleano. Más aún, la localizabilidad es la condición necesaria y suficiente para que dicho retículo sea completo, pues la envolvente de una familia corresponde a su unión a.e.
}
En efecto, sea $b\en\mathbb{R}$ arbitrario:
\begin{itemize}
\item[($\supseteq_{_\text{a.e.}}$)]
Tenemos $A^b_f \setminus (X^b\cap A_f) = A^b_f \setminus X^b \cup A^b_f \setminus A_f = A^b_f \setminus X^b$ (ya que $A^b_f \setminus A_f$). Pero ya probamos que $X^b$ es envolvente de $\mathcal{E}^b$, por lo cual $\mu\bigl(A^b_f \setminus (X^b\cap A_f)\bigr) = \mu\bigl(A^b_f \setminus X^b\bigr) =0$.

\item[($\subseteq_{_\text{a.e.}}$)]
En primera, $(X^b\cap A_f) \setminus A^b_f = X^b\cap A_f \setminus A^b_f$; además, por su definición, 
\[
A_f \setminus A^b_f \;=\; \set{x\en A_f \given f(x) \ge b}% \;\equiv\; \set{f \ge b}
\]
Sea entonces $h\en\mathcal{F}$ arbitraria. Sustituyendo esto último, obtenemos%
\footnote{
    Usamos primero la propiedad elemental: $M\setminus(N\setminus L) = M \cap (L \setminus N)$.
}
\begin{align*}
A^b_h\setminus (A^b_f \setminus A_f)
&\;=\; A^b_h\cap A_f \setminus A^b_f
\;=\; \set*{x\en A_h \given h(x) < b} \cap \set*{x\en A_f \given f(x) \ge b}\\
&\;=\; \set*{x \en A_h \cap A_f \given h(x) < b \leq f(x)}
\;\subseteq\; \set*{x \en A_h \cap A_f \given h(x) \neq f(x)}
\end{align*}
Pero el conjunto del final tiene medida cero ya que $h$ y $f$ coinciden casi en todas partes en $A_h \cap A_f$. Se sigue que $\mu\bigl(A^b_h\setminus (A^b_f \setminus A_f)\bigr) = 0$. Como $h\en\mathcal{F}$ es arbitraria entonces, gracias a que $X^b$ es envolvente de $\mathcal{E}^b$, se tiene $\mu\bigl(X^b\setminus (A^b_f \setminus A_f)\bigr) = 0$. Como $X^b\setminus (A^b_f \setminus A_f) = X^b\cap A_f \setminus A^b_f$ concluimos que $\mu\bigl((X^b\cap A_f) \setminus A^b_f\bigr) = 0$.
\end{itemize}
Hemos probado que $\set{\tilde{g}|_{A_f} < b} \ctpeq \set{f<b}$, donde $b\en\mathbb{R}$ es arbitrario. Solo falta ver que esto implica $\tilde{g}|_{A_f} \ctpeq f$. Esto es así: sea $x \en A_f$, entonces, si $\tilde{g}(x) < f(x)$, por densidad de $\mathbb{Q}$ se sigue que existe $q\en\mathbb{Q}$ tal que $\tilde{g}(x) < q < f(x)$. En otras palabras:
\[
\set{\tilde{g}|_{A_f} < f}
\;=\;\bigcup_{q\en\mathbb{Q}} \set{\tilde{g}|_{A_f} < q} \cap \set{q< f}
\]
pero $\set{\tilde{g}|_{A_f} < q} \cap \set{q< f} \subseteq \set{\tilde{g}|_{A_f} < q} \cap \set{q \le f} = \set{\tilde{g}|_{A_f} < q} \setminus \set{f < q}$; no obstante, este último conjunto tiene medida cero gracias a que $\set{\tilde{g}|_{A_f} < q} \ctpeq \set{f < q}$. Se sigue que $\set{\tilde{g}|_{A_f} < f}$ tiene medida cero al ser unión contable de conjuntos de medida cero.
De forma análoga, 
\[
\set{f<\tilde{g}|_{A_f}}
\;=\;\bigcup_{q\en\mathbb{Q}} \set{f < q}\cap\set{q < \tilde{g}|_{A_f}}
\]
donde $\set{f < q}\cap\set{q < \tilde{g}|_{A_f}} \subseteq \set{f<q} \setminus \set{\tilde{g}|_{A_f} < q}$ tiene medida cero gracias a que tenemos $\set{\tilde{g}|_{A_f} < q} \ctpeq \set{f < q}$, y de esto se sigue que $
\set{f<\tilde{g}|_{A_f}}$ tiene medida cero al ser unión contable de conjuntos de medida cero. Finalmente, esto garantiza que 
\[
\set{\tilde{g}|_{A_f} \neq f}
\;=\; \set{\tilde{g}|_{A_f} < f} \cup \set{f<\tilde{g}|_{A_f}}
\]
tiene medida cero, y así $\tilde{g}|_{A_f} \ctpeq f$.\vspace{1em}\\
\emph{Restrinjamos los valores de $\tilde{g}$ a los reales}:\\
Como $\tilde{g}$ es medible, $\set{\tilde{g} \neq \pm\infty}\equiv\tilde{g}^{-1}(\mathbb{R})\en\mathcal{A}$. Entonces definimos $g\coloneq \tilde{g}\cdot\chi_{\set{\tilde{g}\neq\pm\infty}}$ ($g$ vale cero allá donde $\tilde{g}$ vale $\infty$ o $-\infty$). Tenemos de inmediato que $g \colon X \to \mathbb{R}$ es medible, así que basta probar que sigue siendo extensión c.t.p. de cada $f\en\mathcal{F}$. Esto es inmediato, si $f\en\mathcal{F}$ entonces, ya que $\tilde{g}|_{A_f} \ctpeq f$ y $f$ toma valores reales, tenemos que $\set{\tilde{g}|_{A_f}=\pm\infty} \subseteq \set{\tilde{g}|_{A_f} \neq f}$ tiene medida cero; sin embargo, $\set{\tilde{g}|_{A_f} \neq g|_{A_f}} = \set{\tilde{g}|_{A_f} = \pm \infty}$, por lo cual acabamos de probar que $\set{\tilde{g}|_{A_f} \neq g|_{A_f}}$ tiene medida cero, es decir, $\tilde{g}|_{A_f} \ctpeq g|_{A_f}$. Finalmente, como ya sabíamos que $\tilde{g}|_{A_f} \ctpeq f$, la transitividad de "$\ctpeq$" en funciones medibles (vimos que es rel. de equivalencia para definir $L^{1}(X)$) garantiza $g|_{A_f} \ctpeq f$. Como $f\en\mathcal{F}$ es arbitrario, esto concluye la prueba.
\end{sol}

